\documentclass[11pt]{article}
\usepackage{amsmath}
%\DeclareMathOperator*{\argmax}{argmax}
%\DeclareMathOperator*{\argmin}{argmin} 
\usepackage{amsfonts}
%\usepackage[utf8]{inputenc}
\usepackage[english]{babel}
%\usepackage{amsthm}
%\theoremstyle{definition}
%\newtheorem{definition}{Definition}[section]
%\usepackage{geometry}
\usepackage{mathtools}
%\usepackage{listings}
%
%\usepackage{xcolor}
%
%\definecolor{codegreen}{rgb}{0,0.6,0}
%\definecolor{codegray}{rgb}{0.5,0.5,0.5}
%\definecolor{codepurple}{rgb}{0.58,0,0.82}
%\definecolor{backcolour}{rgb}{0.95,0.95,0.92}
%
%\lstdefinestyle{mystyle}{
%    backgroundcolor=\color{backcolour},   
%    commentstyle=\color{codegreen},
%    keywordstyle=\color{magenta},
%    numberstyle=\tiny\color{codegray},
%    stringstyle=\color{codepurple},
%    basicstyle=\ttfamily\footnotesize,
%    breakatwhitespace=false,         
%    breaklines=true,                 
%    captionpos=b,                    
%    keepspaces=true,                 
%    numbers=left,                    
%    numbersep=5pt,                  
%    showspaces=false,                
%    showstringspaces=false,
%    showtabs=false,                  
%    tabsize=2
%}
%
%\lstset{style=mystyle}
%
%\usepackage{bbm}
%\usepackage{physics}
%\usepackage{parskip}
%\usepackage{hyperref}
%\usepackage[linesnumbered,ruled,vlined]{algorithm2e}
%\DontPrintSemicolon
%\usepackage{varioref}
%\labelformat{algocf}{\textit{Algorithm}\,(#1)}
%
%\usepackage{tikz}
%\usetikzlibrary{shapes.geometric, arrows, calc, fit, positioning}
%\tikzstyle{bluebox} = [rectangle, rounded corners, minimum width=3cm, minimum height=1cm,text centered, draw=black, fill=blue!30]
%\tikzstyle{arrow} = [thick,->,>=stealth]
%
%\usepackage{hyperref}
%\hypersetup{
%    colorlinks,
%    citecolor=black,
%    filecolor=black,
%    linkcolor=black,
%    urlcolor=black
%}

\author{Greg Strabel}
\title{Information Criterion}
\begin{document}
\maketitle

\tableofcontents

\section{Summary}
This document provides a brief overview of Akaike and Bayesian Information Criterion.

\section{Assumptions}

We have a set $\mathbf{X} = \{ x_i \}_{i = 1}^{N}$ of observations from a parametric family of distributions $p(\mathbf{X} \vert \theta)$ where the true value $\theta_0$ is unknown. We assume that $p(\mathbf{X} \vert \theta)$ is twice differentiable in $\theta$ and $\theta \in \mathbb{R}^k$.

\section{Akaike Information Criteria}

Let $\hat{\theta}$ be the maximum likelihood estimator of $\theta$ and define 
\begin{equation}
d(\theta') \coloneqq \left[ \int_{\Omega}-2\ln{p(\mathbf{X} \vert \theta)} d\mathbb{P}(\mathbf{X} \vert \theta) \right] \vert_{\theta=\theta'}
\end{equation}

Then

\begin{align}
\begin{split}
\mathbb{E} \left[ d \left( \hat{\theta} \right) \right] &= \underbrace{\mathbb{E} \left[ -2\ln{p(\mathbf{X} \vert \hat{\theta})} \right]}_{\text{A}} \\
& \quad + \underbrace{ \mathbb{E} \left[ 2 \ln{p( \mathbf{X} \vert \hat{\theta})} - 2 \ln{p( \mathbf{X} \vert \theta_0)} \right] }_{\text{B}} \\
& \quad + \underbrace{ \mathbb{E} \left[ d(\hat{\theta}) + 2 \ln{p(\mathbf{X} \vert \theta_0)} \right] }_{\text{C}}
\end{split}
\end{align}

Using a second order Taylor expansion around $\hat{\theta}$,

\begin{align}
\begin{split}
B &= 2 \mathbb{E} \left[ \ln{p(\mathbf{X} \vert \hat{\theta})} - 
\{ 
\ln{p(\mathbf{X} \vert \hat{\theta})} 
+ \frac{\partial}{\partial \theta} \ln{p(\mathbf{X} \vert \theta)} \vert_{\theta=\hat{\theta}} (\theta_0 - \hat{\theta}) \} \right]  \\
& \quad - 2\mathbb{E} \left[ \frac{1}{2}(\theta_0 - \hat{\theta})' \frac{\partial^2}{\partial\theta\partial\theta'}\ln{p(\mathbf{X} \vert \theta)} \vert_{\theta=\hat{\theta}}(\theta_0 - \hat{\theta}) + o_p(1) \right] \\
&= 2\mathbb{E} \left[ \frac{1}{2}(\theta_0 - \hat{\theta})' \frac{-\partial^2}{\partial\theta\partial\theta'}\ln{p(\mathbf{X} \vert \theta)} \vert_{\theta=\hat{\theta}}(\theta_0 - \hat{\theta}) + o_p(1) \right]
\end{split}
\end{align}

The term $(\theta_0 - \hat{\theta})' \frac{-\partial^2}{\partial\theta\partial\theta'}\ln{p(\mathbf{X} \vert \theta)} \vert_{\theta=\hat{\theta}}(\theta_0 - \hat{\theta})$ converges asymptotically to $\chi^2_k$ by the Central Limit Theorem and the Continuous Mapping Theorem. Therefore $B \rightarrow_p k$.

Likewise, using a second order Taylor expansion around $\theta_0$,

\begin{align}
\begin{split}
C &= \mathbb{E} \left[ d(\theta_0) + \frac{\partial}{\partial\theta}d(\theta) \vert_{\theta=\theta_0}(\hat{\theta} - \theta_0) \right] \\
& \quad + \mathbb{E} \left[ \frac{1}{2}(\hat{\theta} - \theta_0)' \frac{\partial^2}{\partial \theta \partial \theta'}d(\theta) \vert_{\theta=\theta_0}(\hat{\theta} - \theta_0) + o_p(1) + 2\ln{q(\mathbf{X} \vert \theta_0)} \right] \\
&= \mathbb{E} \left[ \frac{1}{2}(\hat{\theta} - \theta_0)' \frac{\partial^2}{\partial \theta \partial \theta'}d(\theta) \vert_{\theta=\theta_0}(\hat{\theta} - \theta_0) + o_p(1) \right] \\
&= \mathbb{E} \left[ (\hat{\theta} - \theta_0)'\frac{-\partial^2}{\partial\theta\partial\theta'}\ln{p(\mathbf{X} \vert \theta)} \vert_{\theta=\theta_0}(\hat{\theta} - \theta_0) + o_p(1) \right]
\end{split}
\end{align}

Because $(\theta_0 - \hat{\theta})' \frac{-\partial^2}{\partial\theta\partial\theta'}\ln{p(\mathbf{X} \vert \theta)} \vert_{\theta=\theta_0}(\theta_0 - \hat{\theta})$ converges asymptotically to $\chi^2_k$, $C \rightarrow_p k$.

It follows that

\begin{equation}
\mathbb{E} \left[ d\left( \hat{\theta} \right) \right] = \mathbb{E} \left[ -2\ln{p(\mathbf{X} \vert \hat{\theta})} \right] + 2k + o(1)
\end{equation}

It follows that the Akaike Information Criterion,

\begin{align}
\begin{split}
AIC &\coloneqq 2k - 2\ln{p(\mathbf{X} \vert \hat{\theta})}
\end{split}
\end{align}

satisfies

\begin{align}
\begin{split}
\mathbb{E} \left[ AIC \right] &= 2k + 2 \mathbb{E} \left[- \ln{p(\mathbf{X} \vert \hat{\theta})} \right] \\
&= \mathbb{E} \left[ d \left( \hat{\theta} \right) \right] + o(1)
\end{split}
\end{align}

\section{Bayesian Information Criteria}

Let $p(\theta)$ be a prior on $\theta$. Then using a second order Taylor expansion around $\hat{\theta}$ and Laplace's Approximation, we have

\begin{align}
\begin{split}
p(\mathbf{X}) &= \int_{\Theta}p(\mathbf{X} \vert \theta) p(\theta)d\theta \\
&= \int_{\Theta}\exp{\left[\ln{p(\mathbf{X} \vert \theta)}\right]}p(\theta)d\theta \\
&= \int_{\Theta}\exp{\left[
\ln{p(\mathbf{X} \vert \hat{\theta)}} + \frac{\partial}{\partial\theta}\ln{p(\mathbf{X} \vert \theta)} \vert_{\theta=\hat{\theta}}(\theta - \hat{\theta})
\right]}p(\theta)d\theta \\
& \quad + \int_{\Theta}\exp{\left[ \frac{1}{2}(\theta - \hat{\theta})'\frac{\partial^2}{\partial\theta\partial\theta'}\ln{p(\mathbf{X} \vert \theta)} \vert_{\theta=\hat{\theta}}(\theta-\hat{\theta}) + o_p(1)
\right]}p(\theta)d\theta \\
& \approx \int_{\Theta}p(\mathbf{X} \vert \hat{\theta})\exp{\left[ \frac{n}{2}(\theta - \hat{\theta})'\frac{\partial^2}{\partial\theta\partial\theta'}\frac{1}{n}\ln{p(\mathbf{X} \vert \theta)} \vert_{\theta=\hat{\theta}}(\theta-\hat{\theta}) \right]} \exp{\left[o_p(1)\right]}p(\hat{\theta})d\theta \\
& \approx p(\mathbf{X} \vert \hat{\theta}) \left( \frac{2\pi}{n} \right)^{k/2}\left[\det{\left(\frac{-\partial^2}{\partial\theta\partial\theta'}\frac{1}{n}\ln{p(\mathbf{X} \vert \theta)} \vert_{\theta=\hat{\theta}}\right)}\right]^{-1/2}p(\hat{\theta})
\end{split}
\end{align}

It follows that

\begin{align}
\begin{split}
\ln{p(\mathbf{X})} & \approx \ln{p(\mathbf{X} \vert \hat{\theta})} + \frac{k}{2}\ln{2\pi} - \frac{k}{2}\ln{n} \\
& \quad - \frac{1}{2}\ln{\left[\det{\left(\frac{-\partial^2}{\partial\theta\partial\theta'}\frac{1}{n}\ln{p(\mathbf{X} \vert \theta)} \vert_{\theta=\hat{\theta}}\right)}\right]} + \ln{p(\hat{\theta})} \\
&= \ln{p(\mathbf{X} \vert \hat{\theta})} - \frac{k}{2}\ln{n} + O_p(1)
\end{split}
\end{align}

so that

\begin{align}
\begin{split}
BIC & \coloneqq k\ln{n} - 2\ln{p(\mathbf{X} \vert \hat{\theta})} \\
& \approx -2\ln{p(\mathbf{X})} + O_p(1)
\end{split}
\end{align}

%\bibliographystyle{alpha} % We choose the "plain" reference style
%\bibliography{information_criterion} % Entries are in the refs.bib file

\end{document}